% A.68 candidate: notes only, unregistered pending independent audits.
\subsection{The two nonstable dihedral genus-seven covers}

\begin{statement}{A.68}{Open}{Candidate: the nonstable dihedral genus-seven covers have cusp fixed groups}
Let
\[
 C_{\pm}:v^2=w^{16}\pm6w^8+1,
 \qquad S=\{P_{0,+},P_{0,-},P_{\infty,+},P_{\infty,-}\},
\]
where the signs distinguish the values of \(v\) at zero and of
\(v/w^8\) at infinity. Equip these punctured curves with the actual
degree-sixteen covers of \(E\setminus\{O\}\) described below.
For every coefficient prime, every induced cusp packet, and every
open derived Galois subgroup preserving the selected cover and cusp,
the fixed geometric pro-\(\ell\) subgroup is exactly the selected
cusp inertia. Their procyclic normalizer assertions also hold, by
the nilpotent-closure theorem A.59.

The new geometric input is that their proper Jacobians have potential
good reduction of \(2\)-rank zero at two and that every cusp difference
is annihilated by eight. This is a result for these actual covers;
it makes no claim about arbitrary higher dihedral covers.
\end{statement}

\begin{proof}
We first work with the minus-six member \(C\). The substitution
\(w\mapsto\zeta_{16}w\), after a finite field extension, identifies
the two members and their four cusp sets.

\paragraph{The actual cover and its quotient maps.}
The degree-sixteen map to the original elliptic curve is
\[
 X=\frac{(w^8-1)^2}{4w^8},\qquad
 Y=\frac{v(w^{16}-1)}{8w^{12}},\qquad Y^2=X^3-X.
                                                               \tag{A.68.1}
\]
It also factors as elliptic duplication after the actual degree-four map
\[
 x=\frac{w^8-1+v}{2},\qquad y=w^4(x-1),
 \qquad y^2=x^3-x,\qquad w^4=\frac{y}{x-1}.           \tag{A.68.2}
\]
Substitution and the elliptic duplication formulas verify
\([2](x,y)=(X,Y)\). The divisor of \(y/(x-1)\) on the elliptic
curve is
\[
                 T_0+T_{-1}-T_1-O.
\]
Consequently (A.68.2) is totally ramified of index four at precisely
the four cusps. Duplication is unramified, so (A.68.1) is finite
\'{e}tale away from these cusps. Its deck group is the actual dihedral
group of order sixteen generated by
\[
 r(w,v)=(\zeta_8w,-v),\qquad s(w,v)=(1/w,-v/w^8).
\]
Indeed both preserve (A.68.1), satisfy \(r^8=s^2=1\) and
\(srs=r^{-1}\), and give sixteen distinct automorphisms.
The hyperelliptic equation has sixteen distinct roots, so \(g(C)=7\).

Besides the original deck action, put
\(\mu(w,v)=(\zeta_8w,v)\). The quotient by \(\mu^4\) is the
actual genus-three curve
\[
 C_4:\ v^2=u^8-6u^4+1,\qquad u=w^2.                 \tag{A.68.3}
\]
There is also an actual map to the genus-two curve
\[
 \begin{split}
 B:&\quad W^2=(T+2)(T^4-4T^2-4),\\
 &T=w^2+w^{-2},\qquad W=(w+w^{-1})v/w^4.
 \end{split}                                                   \tag{A.68.4}
\]
The polynomial on the right has discriminant \(-2^{20}\),
so its smooth proper model has genus two. To see the quotient
construction explicitly, first quotient by
\(h(w,v)=(1/w,v/w^8)\), obtaining
\[
 t=w+w^{-1},\quad z=v/w^4,\qquad
 z^2=t^8-8t^6+20t^4-16t^2-4.
\]
Then take invariants \(u=t^2,W=tz\) under \((t,z)\mapsto(-t,-z)\)
and set \(T=u-2\). These substitutions prove (A.68.4), including
the function-field identification rather than merely its branch data.

\paragraph{The full Jacobian decomposition comes from those maps.}
Use the basis
\(\omega_j=w^{j-1}dw/v\), \(1\le j\le7\), of
\(H^0(C,\Omega^1)\). The pullbacks from \(C_4\) span
\(\omega_2,\omega_4,\omega_6\). For (A.68.4), direct differentiation gives
\[
 \frac{dT}{W}=2(\omega_5-\omega_3),\qquad
 \frac{T\,dT}{W}=2(\omega_7-\omega_5+\omega_3-\omega_1).
                                                               \tag{A.68.5}
\]
Thus its pullback plane is spanned by
\(\omega_1-\omega_7\) and \(\omega_3-\omega_5\).
Since \(\mu^*\omega_j=\zeta_8^j\omega_j\), this plane and its
\(\mu^*\)-translate are complementary in the four-dimensional
odd-index space: for each pair \((j,8-j)=(1,7),(3,5)\),
the two eigenvalues are distinct. The three actual maps therefore
induce an isomorphism on rational cotangent spaces of their product
Jacobians, hence an isogeny
\[
                       J_C\sim J_{C_4}\times J_B^2.            \tag{A.68.6}
\]
All maps and the isogeny are defined after a finite number-field
extension; this is sufficient for potential reduction and the
derived-group formulation in A.41--A.42.

\paragraph{Exact Igusa invariants of the new genus-two factor.}
We use the standard Igusa invariants \(J_{2i}\), with
\(J_{10}=2^{-12}\operatorname{disc}(f)\) for our monic quintic.
For clarity the intermediate Igusa--Clebsch invariants are denoted
\(I^{\rm IC}_{2i}\), distinct from Liu's auxiliary invariants below.
Homogenize the quintic in (A.68.4) to a binary sextic \(F\).
With the normalized transvectant
\[
 (f,g)_k=\frac{(m-k)!(n-k)!}{m!n!}
 \sum_{j=0}^k(-1)^j\binom{k}{j}
 \frac{\partial^kf}{\partial x^{k-j}\partial z^j}
 \frac{\partial^kg}{\partial x^j\partial z^{k-j}},
                                                               \tag{A.68.7}
\]
put \(i=(F,F)_4\), \(d=(i,i)_2\), and
\(A=(F,F)_6,B_0=(i,i)_4,C_0=(i,d)_4\). Exact polynomial
differentiation gives
\[
 A=-8/5,\qquad B_0=1792/1875,\qquad C_0=-149504/46875.
\]
The standard conversion
\[
 \begin{split}
 I^{\rm IC}_2&=-120A,\\
 I^{\rm IC}_4&=-720A^2+6750B_0,\\
 I^{\rm IC}_6&=8640A^3-108000AB_0+202500C_0
 \end{split}
\]
gives \((192,4608,-516096)\), while
\(I^{\rm IC}_{10}=\operatorname{disc}(f)=-1048576\).
Using the conversions in D\k{a}browski--Sadek, equation (3),
\[
 \begin{aligned}
 J_2&=I^{\rm IC}_2/8,&
 J_4&=(4J_2^2-I^{\rm IC}_4)/96,\\
 J_6&=(8J_2^3-160J_2J_4-I^{\rm IC}_6)/576,&
 J_8&=(J_2J_6-J_4^2)/4,
 \end{aligned}
\]
we obtain
\[
                 (J_2,J_4,J_6,J_8,J_{10})
                    =(24,-24,1248,7344,-256).                 \tag{A.68.8}
\]
The exact certificate also reproduces the official PARI example
\(x^5+3x^2-4\), after the explicitly checked weighted rescaling from \(P\) to \(4P\).
Every criterion below is a quotient of equal weights, so this
independent normalization check rules out a scaling discrepancy.

\paragraph{Liu's residue-two criterion determines the stable fibre.}
Liu, \emph{Courbes stables de genre 2 et leur sch\'ema de modules},
Theorem 1 on printed page 204, applies in residue characteristic two
with \(\varepsilon=4\) and \(I_{2\varepsilon}=J_8\). His auxiliary
invariants, defined on printed page 203, are
\[
 \mathcal I_4=J_2^2-24J_4=1152,
 \quad
 \mathcal I_{12}=-8J_4^3+9J_2J_4J_6-27J_6^2-J_2^2J_8
                 =-52641792.
\]
Their valuations, with \(v_2(2)=1\), are
\(v_2(\mathcal I_4)=7\), \(v_2(\mathcal I_{12})=14\),
\(v_2(J_8)=4\), and \(v_2(J_{10})=8\). Therefore the three
expressions in Theorem 1\((V_*)\) have valuations
\[
 \begin{array}{c|ccc}
 \text{expression}&\mathcal I_4^4/J_8^2&J_{10}^4/J_8^5&
                              \mathcal I_{12}^4/J_8^6\\\hline
 v_2&20&12&32.
 \end{array}                                                   \tag{A.68.9}
\]
They lie in the maximal ideal, so the stable fibre has two
irreducible components meeting in one point. The two additional
expressions in part \((V)\) have valuations
\[
 v_2\!\left(\frac{\mathcal I_4^{12}}{J_{10}^4J_8}\right)=48,
 \qquad
 v_2\!\left(\frac{\mathcal I_{12}^{4}}{J_{10}^4J_8}\right)=20.
                                                               \tag{A.68.10}
\]
Both are integral, so the two components are smooth elliptic curves.
Moreover the theorem's formulas are
\[
 (j_1j_2)^4=
  \overline{\mathcal I_4^{12}/(J_{10}^4J_8)},\qquad
 (j_1+j_2)^4=2^6\cdot3^3+
  \overline{\mathcal I_{12}^{4}/(J_{10}^4J_8)}.
\]
Both right-hand sides vanish in characteristic two. Hence
\(j_1=j_2=0\), and both elliptic components are supersingular.
Thus \(J_B\) has potential good reduction with \(2\)-rank zero.

The factor \(J_{C_4}\) has this same property by the proper-factor
calculation in A.37: its elliptic factors have integral \(j\)-invariants
reducing to zero at two. Potential good reduction is preserved by
isogeny (equivalently use an auxiliary prime in the
N\'eron--Ogg--Shafarevich criterion), and the \(2\)-rank of good
abelian reduction is isogeny invariant. Equation (A.68.6) therefore
proves potential good reduction of \(J_C\) with \(2\)-rank zero.
In particular the stable proper fibre of \(C\) is of compact type.

\paragraph{All four cusp differences are actual eight-torsion.}
The following identities hold in the function field of \(C\):
\[
 \begin{split}
 (v-w^8+1)(v+w^8-1)&=-4w^8,\\
 (v-w^8-1)(v+w^8+1)&=-8w^8.
 \end{split}                                                   \tag{A.68.11}
\]
There are no zeros or poles outside the four cusps. Expanding
\(\sqrt{1-6u^8+u^{16}}=1-3u^8-4u^{16}+O(u^{24})\) at zero,
and using \(u=w^{-1}\) at infinity, determines the exact divisors:
\[
 \begin{split}
 \operatorname{div}(v-w^8+1)&=8P_{0,-}-8P_{\infty,-},\\
 \operatorname{div}(v+w^8-1)&=8P_{0,+}-8P_{\infty,+},\\
 \operatorname{div}(v-w^8-1)&=8P_{0,+}-8P_{\infty,-},\\
 \operatorname{div}(v+w^8+1)&=8P_{0,-}-8P_{\infty,+}.
 \end{split}                                                   \tag{A.68.12}
\]
These four relations imply \(8[P-Q]=0\) for any two cusps.
After a finite extension giving the good Jacobian and separated
marked stable curve, each such divisor specializes to zero because
the special Jacobian has no nonzero geometric \(2\)-power torsion.
The torsion-divisor criterion proved in A.41 then says that the
joining path between any two cusp components contains only rational
components. Every cusp has a distinct companion cusp. A.42 applies
to every coefficient prime, every cusp packet and every open derived
subgroup, proving the stated full fixed-subgroup equality. A.59
independently supplies the normalizer assertion for the actual
dihedral closure (A.68.1).

All finite identities used here are checked in the retained program
\path{notes/fgh_C8_compact_type_certificate.py}; the general
fixed-group implication is the proved A.41--A.42 theorem, rather than
an inference from a finite computation.
\end{proof}
